\subsection{Experiment~14: Where the Sparse Code Is Formed}
\label{sec:exp14-residue-sparse}

The encoder deployed in this system projects a protein's pooled embedding into a
sparse, ontology-aligned code, and retrieval then runs on that code rather than on
the dense vector. This experiment asks a question the deployed design never puts:
whether the sparse selection should happen \emph{before} the residues are pooled
rather than after.

The question is not cosmetic, because selection and averaging do not commute. Taking
the largest entries of an average is a different operation from averaging the largest
entries of each term, and the two disagree whenever an atom is a narrow runner-up in
many residues and a decisive winner in few. Pooling first lets such an atom bank
every second place it never won; selecting first never admits a second place at all.
Once a pooled vector exists the distinction is unrecoverable, since which residue
activated which atom has already been averaged away, so the two orders cannot be
compared by any post-hoc transformation of the stored embeddings. They can only be
compared by running the language model again.

\subsubsection{Design}

All encoders are fitted on the same proteins, with the same objective, budget and
random seed, and scored on the same held-out proteins that no encoder has seen. The
population being scored is the temporal evaluation delta; the proteins that the
deployed encoder was itself fitted on are removed before anything is measured, since
scoring it on those would credit it with an advantage that has nothing to do with the
mechanism under test.

Two references are reported for every arm, because there are two questions with
different answers. Against a \emph{pooled control}, an encoder fitted on the same
proteins with the same objective and budget, the comparison isolates the mechanism:
whatever an arm wins there is not explained by having seen more data or a different
loss. Against the \emph{deployed encoder}, frozen, the comparison asks whether any of
it improves on what the system currently runs, a coarser question in which the
training budget is part of the comparison rather than a confound to be removed.

The primary screen is the mean functional similarity between a protein and its ten
nearest neighbours by code cosine, over ancestor closures propagated through the
ontology. It is used because it is what nearest-neighbour transfer consumes, and its
validity as a proxy is tested rather than assumed in
Section~\ref{sec:exp14-fmax}. Every difference carries a bootstrap interval paired on
proteins, and every result is reported by protein length band.

\subsubsection{A resolution floor, measured rather than assumed}
\label{sec:exp14-floor}

Before any arm is ranked, the study's own reproducibility is measured. The same
recipe, the same seed, fitted twice on two extractions verified identical to the
last bit, produces encoders that differ by $0.0013$ in mean purity, with a paired
interval that excludes zero.

That number is not noise in the estimator; it is variance in the fit. A bootstrap
paired on proteins resamples the population and is silent about the optimiser, and
two mechanisms feed the gap: the order in which minibatches arrive, and the fact that
a selection of the largest entries is unstable under it. A shift of one part in ten
million swaps a near-tied atom, and the code then differs by that atom's whole
magnitude. An independent estimate agrees: two aggregation routes that are
algebraically identical for proteins occupying a single window score $0.0006$ apart,
again with a separating interval.

Every difference below roughly $0.0015$ is therefore reported as unresolved
regardless of its interval. Three of the intervals in this section exclude zero and
are nonetheless declared inconclusive on that basis. Stating the floor first, rather
than discovering it when a favourable difference needs defending, is what makes the
ranking that follows meaningful.

\subsubsection{The mechanism}

Table~\ref{tab:exp14-mechanism} gives the ordering. Selecting atoms per residue
before pooling is worth more than three hundredths of purity against the pooled
control, roughly a quarter of the control's own absolute score, and the advantage
survives at every neighbourhood size from one to fifty and grows with it.

\begin{table}[htbp]
  \centering
  \caption{Encoding mechanisms against a pooled control fitted on the same proteins
  with the same objective and budget. Held out on 2{,}646 proteins.
  $^{\dagger}$~below the resolution floor of Section~\ref{sec:exp14-floor}.}
  \label{tab:exp14-mechanism}
  \begin{tabular}{lrrl}
\toprule
    \textbf{Encoding} & \textbf{Purity} & \textbf{$\Delta$} & \textbf{95\,\% interval} \\
\midrule
    Per residue, 4 atoms, all pairs & 0.2626 & $+0.0414$ & $[+0.0390, +0.0438]$ \\
    Per residue, 4 atoms & 0.2570 & $+0.0358$ & $[+0.0336, +0.0381]$ \\
    Per residue, 16 atoms & 0.2556 & $+0.0344$ & $[+0.0321, +0.0367]$ \\
    Per residue, 32 atoms & 0.2535 & $+0.0323$ & $[+0.0300, +0.0345]$ \\
    Per residue, 128 atoms & 0.2465 & $+0.0253$ & $[+0.0234, +0.0273]$ \\
    Per residue, no selection & 0.2284 & $+0.0071$ & $[+0.0059, +0.0084]$ \\
    Deployed encoder (pooled) & 0.2284 & $+0.0071$ & $[+0.0058, +0.0085]$ \\
    Pooled vector, no sparse code & 0.2180 & $-0.0033$ & $[-0.0049, -0.0017]$ \\
\bottomrule
\end{tabular}

\end{table}

The number of atoms retained per residue behaves as a plateau rather than an optimum.
Moving from 128 to 32 is worth five times the resolution floor; from 32 to 16, less
than twice it; and the steps from 16 to 8 and from 8 to 4 fall below it. Since four
atoms per residue is not measurably worse than sixteen and is thirty-two times
cheaper than the deployed 128, the choice inside the plateau is made on cost and not
on signal. This is the rare case where the recommended setting is defended by an
absence of difference, which is only admissible because the floor was measured first.

\subsubsection{Interventions that did not help}

Table~\ref{tab:exp14-negatives} collects the changes that were expected to help and
did not. They are reported because each one closes a direction, and because two of
them contradict standard practice.

\begin{table}[htbp]
  \centering
  \caption{Interventions measured against the identical arm without them.
  $^{\dagger}$~below the resolution floor.}
  \label{tab:exp14-negatives}
  \begin{tabular}{lrl}
\toprule
    \textbf{Intervention} & \textbf{$\Delta$} & \textbf{Against} \\
\midrule
    Binary residue values & $-0.0615$ & the same without it \\
    Ranking loss & $-0.0118$ & the same without it \\
    Hard negatives & $-0.0045$ & the same without it \\
    Aggregate by window & $-0.0031$ & the same without it \\
    All four layers & $-0.0047$ & the same without it \\
    More pairs at fixed cost & $-0.0032$ & the same without it \\
    Four times the corpus & $-0.0002$$^{\dagger}$ & the same without it \\
    Pretrain on the large corpus & $+0.0001$$^{\dagger}$ & the same without it \\
\bottomrule
\end{tabular}

\end{table}

Hard negative mining, which the deployed encoder uses, costs accuracy here. The
mechanism is specific rather than mysterious: the objective is a regression that asks
the code cosine to equal a functional similarity, and the calibration of that scale
is carried by the mass of easy pairs. Oversampling the pairs the code currently gets
wrong biases the estimate toward suppressing cosines and distorts the scale it is
trying to fit. A ranking loss, tried as the natural alternative, was worse still, and
because it was worse for signed codes as well as for binary ones the failure is
attributable to the objective and not to an interaction with the code. The lesson is
that the objective and the sampling scheme are one decision and cannot be chosen
apart.

Binary residue values fail for a reason that is geometric and was confirmed by direct
measurement rather than inferred from the score. A code whose entries cannot be
negative has cosines confined to a narrow, high band, while the functional
similarities it is asked to reproduce are concentrated near zero. It is not that the
representation is imprecise; it is that it cannot express dissimilarity at all. The
practical consequence is that a quantised code must keep its sign, and that a
representation of this kind cannot be trained by ordinary backpropagation at all,
since replacing a value by a membership indicator leaves the output independent of
the weights and the backward pass finds no path.

Aggregating at the window rather than the residue also loses, and it loses most
clearly where it can act: on the proteins long enough to occupy more than one window
it is five times worse than its diluted figure over the whole population. Giving each
window equal weight discards the fact that windows hold different numbers of
residues, which per-residue pooling already respected.

Finally, the size of the training corpus saturates early. Four thousand proteins to
fifteen thousand is worth about twice the floor; fifteen thousand to sixty thousand
is indistinguishable from zero, and pre-training on the larger corpus and adapting
afterwards is indistinguishable from not using it at all. What does matter, by a
factor of six, is whether the training proteins are drawn from the population being
scored. The lever that the literature would reach for first is the one that is
already exhausted, and the lever that is not usually described as a lever carries the
effect.

\subsubsection{Do the layers of the language model carry different signal?}

A comparison across layers is impossible on raw activations, and the reason is worth
recording. The middle layers of the backbone used here carry activations four orders
of magnitude larger than the last: typical vector norms run from under two at the
output to nearly one hundred thousand in the middle. A map initialised and optimised
identically across such a range is not in the same regime for each layer, so a first
comparison ranked the input scale and not the layer, and it silently broke the
dictionary comparison as well, since a concatenation of four layers at those relative
scales simply is the two largest of them.

After standardising each layer with statistics fitted on the training split alone,
the ordering is monotone in depth and the last layer wins. A learned mixture over
layers, initialised at an equal average and free to move, converges to weights that
increase monotonically with depth and places nearly half its mass on the final layer;
no combination of layers beats that layer alone, and the best of them merely ties it
within the floor. The practical consequence is that extracting further layers is not
worth its cost, and that a single layer is both the best and the cheapest input.

This disagrees with an earlier measurement in this work, which placed the middle
layers ahead. That measurement was made on pooled embeddings and a different task,
and the two are not the same question. The disagreement is recorded rather than
resolved, because resolving it would require rerunning the earlier protocol under the
present one and no such run was made.

\subsubsection{From the screen to the metric}
\label{sec:exp14-fmax}

Every result above is measured on a functional screen. The screen was chosen because
it is what nearest-neighbour transfer consumes, but the claim being made is about
annotation quality, and no amount of purity is evidence for it until the two are put
side by side.

The full task was therefore run: a reference bank of sixty thousand annotated
proteins, the held-out queries retrieving their nearest neighbours in it, terms
transferred with a score proportional to the similarity mass supporting them, and the
result scored by the maximum F-measure per aspect. The retrieval and the transfer are
identical across arms; only the encoder differs.

\begin{table}[htbp]
  \centering
  \caption{Maximum F-measure on the full retrieval and transfer task, against the
  deployed encoder. Intervals are 2{,}000 bootstrap draws paired on proteins, with
  the measure recomputed on each draw.}
  \label{tab:exp14-fmax}
  \begin{tabular}{lrrrl}
\toprule
    \textbf{Aspect} & \textbf{Proposed} & \textbf{Deployed} & \textbf{$\Delta$} & \textbf{95\,\% interval} \\
\midrule
    MFO & 0.7192 & 0.7126 & $\mathbf{+0.0065}$ & $[+0.0016, +0.0113]$ \\
    BPO & 0.5293 & 0.5079 & $\mathbf{+0.0214}$ & $[+0.0170, +0.0255]$ \\
    CCO & 0.6897 & 0.6798 & $\mathbf{+0.0100}$ & $[+0.0060, +0.0140]$ \\
\bottomrule
\end{tabular}

\end{table}

All three aspects improve and every interval excludes zero. The largest gain falls on
biological process, which is the aspect this thesis has characterised as the
frontier, and it is more than three times the gain on molecular function.

Two features of the comparison make the result stronger rather than weaker. The
reference bank is the corpus the deployed encoder was fitted on, so it encodes that
bank in the domain it was trained for, while the proposed encoder encodes it out of
domain. And the ordering the screen produced is the ordering the task produces, which
retroactively justifies refining eleven axes on a screen that costs minutes instead
of on a metric that costs hours. That the proxy held is a methodological result in its
own right, and it is one that could only be claimed after the fact.

\subsubsection{Length is not a nuisance parameter}

\begin{table}[htbp]
  \centering
  \caption{Change in maximum F-measure against the deployed encoder, by protein
  length in residues. $^{*}$~the 95\,\% interval excludes zero.}
  \label{tab:exp14-fmax-bands}
  \begin{tabular}{lrrrr}
\toprule
    \textbf{Aspect} & \textbf{$\leq$512} & \textbf{512-1024} & \textbf{1024-2048} & \textbf{$>$2048} \\
\midrule
    MFO & $+0.0035$ & $+0.0066$ & $+0.0165$$^{*}$ & $+0.0327$$^{*}$ \\
    BPO & $+0.0242$$^{*}$ & $+0.0187$$^{*}$ & $+0.0118$$^{*}$ & $+0.0108$ \\
    CCO & $+0.0127$$^{*}$ & $+0.0068$ & $+0.0006$ & $-0.0006$ \\
\midrule
    $n$ & 1481 & 799 & 285 & 81 \\
\bottomrule
\end{tabular}

\end{table}

Table~\ref{tab:exp14-fmax-bands} is the reason this thesis reports nothing in
aggregate. The three aspects have opposite length profiles. On molecular function the
advantage is absent in short proteins and grows monotonically with length until it is
five times its own average. On biological process it is largest in short proteins and
fades. On cellular component it exists only in short proteins and is indistinguishable
from zero beyond a thousand residues.

A single pooled figure would have reported a flat improvement three times over and
would have been wrong about the mechanism in all three cases. The profiles also point
somewhere: the aspect whose advantage grows with length is the one whose evidence is
most likely to be local to a domain rather than distributed over the chain, which is
exactly what a code that selects per residue should be able to keep and a pooled
average should be expected to lose.

\subsubsection{How many bits a stored value needs}
\label{sec:exp14-bits}

A code that is better is only deployable if it is also affordable to store and to
scan. The retrieval format holds a code as pairs of an atom index and its value, so a
dictionary of this width costs eleven bits to name an atom and the value precision is
what remains negotiable.

Two ways of reducing it were compared. Quantising the code after training is free to
implement and loses whatever it loses. Training with the quantiser inside the forward
pass, with the gradient routed straight through it, is the standard remedy and is
expected to recover most of that loss.

It does the opposite. Training with the quantiser is worse at four bits, at two and
at one, and by an increasing margin as the levels get coarser. Two mechanisms are
plausible and neither was verified: the straight-through gradient ignores that the
scale is itself a function of the code being quantised, so the objective moves with
the map; and the estimator here must pass through two stacked discrete operations
rather than one, since a selection of the largest entries already sits below the
quantiser. The finding is reported as an empirical negative, because a plausible
mechanism for a result is not the same as having measured it.

Applied after training, the picture is favourable and aspect-dependent
(Table~\ref{tab:exp14-quantisation}). Four bits is indistinguishable from half
precision on all three aspects while removing forty-four per cent of the bytes, and
it still improves on the deployed encoder on all three with separating intervals. Two
bits removes half the bytes and costs nothing measurable on biological process, which
is the aspect this thesis cares most about, but it gives back the whole gain on
molecular function and lands in a tie with the deployed system there. Four bits is
therefore the operating point, and two bits is available when storage is the binding
constraint and that trade is acceptable.

\begin{table}[htbp]
  \centering
  \caption{Maximum F-measure against the bits kept per stored value, with the bytes a
  code occupies in the sparse retrieval format. $^{*}$~the 95\,\% interval excludes
  zero.}
  \label{tab:exp14-quantisation}
  \begin{tabular}{lrrrr}
\toprule
    \textbf{Value precision} & \textbf{Bytes} & \textbf{MFO} & \textbf{BPO} & \textbf{CCO} \\
\midrule
    Half precision & 432 & 0.7192 & 0.5293 & 0.6897 \\
    Four bits & 240 & 0.7211 & 0.5289 & 0.6910 \\
    Two bits & 208 & 0.7122 & 0.5279 & 0.6865 \\
    Deployed encoder & 432 & 0.7126 & 0.5079 & 0.6798 \\
\midrule
    \quad Four bits against half precision & & $+0.0019$ & $-0.0005$ & $+0.0013$ \\
    \quad Two bits against half precision & & $-0.0070$$^{*}$ & $-0.0014$ & $-0.0032$$^{*}$ \\
\bottomrule
\end{tabular}

\end{table}

One accounting caveat belongs with the byte figures. They describe the sparse format
that a retrieval scan reads, where a code is a list of indices and values. Stored as a
dense vector in the database column, the same code occupies its full width regardless
of how few atoms are non-zero and regardless of their precision, so this saving is one
on the served bank and on scan bandwidth, not on the table.

This also settles, for a second time and by a cleaner contrast, what went wrong with
the binary residue values of Table~\ref{tab:exp14-negatives}. A signed one-bit code
scores far above the non-negative binary one at the same precision. The obstruction
was never how few levels there are; it was that a code confined to non-negative
values cannot express dissimilarity. Two remedies proposed for that diagnosis, a
ranking objective and quantisation-aware training, were both measured and both failed.
A correct account of a mechanism does not license the first cure that follows from it.

\subsubsection{What this establishes, and what it does not}

The mechanism is worth more than the data. At a matched corpus, matched size, and
with the deployed encoder additionally enjoying a mining scheme the proposed one does
not use, moving the selection to the residue level is worth roughly four times what
multiplying the training corpus by fifteen is worth. That ratio, rather than any
single score, is the finding.

Three limitations bound it. The evaluation population is one temporal window and one
reference bank, so the result is a measurement of this environment and not a general
claim about protein language models. The screen and the metric agree here, which is
evidence that the proxy was adequate for this ordering and not that it is adequate in
general. And the comparison holds the backbone fixed throughout; nothing here says
whether a different backbone would move the same axes in the same directions, and the
one axis that was varied within the backbone, the choice of layer, was enough to show
that such questions can invert when a confound of scale is removed.
